\documentclass[11pt]{article}
\usepackage[T1]{fontenc}
\usepackage[utf8]{inputenc}
\usepackage{lmodern}
\usepackage{amsmath,amssymb}
\usepackage{longtable,booktabs,array}
% Compatibility with Pandoc's captionless tables on older longtable versions.
\ifcsname c@none\endcsname\else\newcounter{none}\fi
\usepackage{graphicx}
\usepackage[margin=25mm]{geometry}
\usepackage{hyperref}
\hypersetup{colorlinks=true,linkcolor=blue,urlcolor=blue}
\providecommand{\tightlist}{\setlength{\itemsep}{0pt}\setlength{\parskip}{0pt}}
\providecommand{\pandocbounded}[1]{#1}
\setlength{\emergencystretch}{3em}
\setcounter{secnumdepth}{-1}
\begin{document}
\section{Complex analytic spaces and
analytification}\label{complex-analytic-spaces-and-analytification}

\emph{Written and self-checked by GPT-6.1 Sol (OpenAI) in Codex, Ultra
setting, October 2026. Analytic prerequisite integration by GPT-6 Astra
(OpenAI), Codex, Ultra. The linked analytic bridge was written and
self-checked by Claude Opus 5.5; a full independent review is not
claimed. Independently authored material is CC0; the
preparation-and-division component below retains Demailly's OpenContent
permission.}

An equation with complex coefficients defines both a scheme and a space
of complex points. The second space has many more open sets and
functions: small balls are open, and convergent power series are
allowed. Nevertheless, at a complex point the two local rings have
exactly the same formal completion. This is the local reason that
passing from algebraic coherent sheaves to analytic coherent sheaves
preserves exact sequences. The next lesson will explain why projectivity
turns this local comparison into a global equivalence.

We use the reduced convention: a complex algebraic variety is a reduced
separated scheme of finite type over \(\mathbf C\), and analytic spaces
have reduced structure sheaves. Coherent modules may have torsion,
including modules annihilated by powers of ideals. Reducedness of the
ambient space does not require reduced supports for its modules.

\subsection{1. Analytic equations and their
sheaves}\label{analytic-equations-and-their-sheaves}

For an open set \(W\subset\mathbf C^n\), let \(\mathcal O_W\) be the
sheaf of holomorphic functions. Iterated Cauchy integrals on polydiscs
give Taylor expansions; see
\href{../complex-analytic-spaces-and-coherent-sheaves/holomorphic-functions-of-several-variables.html\#1-polydiscs-and-the-definition-of-holomorphy}{Holomorphic
functions of several variables, Theorems 1.2 and 2.1} for the integral
and coefficient proofs. Demailly's freely available \emph{Complex
Analytic and Differential Geometry} is additional reading. At the origin
its stalk is

\[
H_n=\mathbf C\{z_1,\ldots,z_n\},
\]

the ring of power series converging in some neighbourhood of zero. Its
maximal ideal is \((z_1,\ldots,z_n)\). A germ with nonzero constant term
has a convergent reciprocal after shrinking the neighbourhood; every
other germ belongs to that ideal.

An \textbf{analytic subset} \(A\subset W\) is locally the common zero
set of finitely many holomorphic functions. Its vanishing ideal sheaf
\(\mathcal I_A\) consists of all germs vanishing on \(A\), and its
reduced structure sheaf is

\[
\mathcal O_A=(\mathcal O_W/\mathcal I_A)|_A.
\]

It embeds in the sheaf of continuous complex-valued functions on \(A\).
A reduced analytic space is a Hausdorff locally ringed space locally
isomorphic to one of these models. A map is holomorphic if, in local
embeddings, its coordinate functions are restrictions of holomorphic
functions. This definition is independent of embeddings: a morphism of
local models pulls back the coordinate germs, their representatives
define a map to the ambient affine space, and the defining equations
force its image into the target model.

A module is \textbf{coherent} if it is locally finitely generated and
the relations among every finite family of sections are locally finitely
generated. Finite generation at one stalk is weaker than coherence on a
neighbourhood. That distinction is the content of Oka's theorem.

\textbf{Analytic foundation and reading order.} Preparation and division
are proved in Section 2. After that argument, read the following local
analytic results in \emph{Complex analytic spaces and coherent sheaves},
then return to Section 3. Their proofs concern analytic germs and
coherent sheaves, not GAGA, so none presupposes the comparison theorem
that this course will prove. Jean-Pierre Demailly's freely available
\href{https://www-fourier.univ-grenoble-alpes.fr/~demailly/manuscripts/agbook.pdf}{\emph{Complex
Analytic and Differential Geometry}}, together with the human sources
credited in the bridge, supplies further reading. The
\href{prerequisites.html}{prerequisite and proof guide} records the
exact scope and the remaining dependencies.

\begin{itemize}
\tightlist
\item
  \hyperref[2-local-analytic-algebra]{Preparation and division, Section
  2 below}: if a germ is regular of order \(r\) in the last variable, it
  is a unit times a monic degree-\(r\) polynomial in that variable whose
  other coefficients vanish at the origin. Division has a unique
  polynomial remainder of degree less than \(r\).
\item
  \href{../complex-analytic-spaces-and-coherent-sheaves/coherent-sheaves-and-okas-theorem.html\#2-oka-s-coherence-theorem}{Coherence
  of holomorphic functions, Theorem 2.1}: \(\mathcal O_W\) is coherent.
  Oka's theorem uses division to reduce analytic relations to finitely
  many polynomial coefficient relations in one fewer variable, and
  proves that the generators work near the point as well as at it.
\item
  \href{../complex-analytic-spaces-and-coherent-sheaves/analytic-germs-local-parametrization-and-the-nullstellensatz.html\#5-the-nullstellensatz-and-dimension}{The
  local analytic Nullstellensatz, Theorem 5.1}: for an ideal
  \(J\subset H_n\), the ideal of germs vanishing on its zero germ is
  \(\sqrt J\).
\item
  \href{../complex-analytic-spaces-and-coherent-sheaves/cartans-coherence-theorem-and-complex-spaces.html\#1-the-ideal-sheaf-of-an-analytic-set}{Coherence
  of analytic ideal sheaves, Theorem 1.1}: \(\mathcal I_A\) is coherent
  for every analytic subset. Cartan's theorem constructs uniform local
  equations using finite projections and coefficient interpolation; the
  statement includes singular analytic subsets.
\end{itemize}

Using these coherence results, coherent modules on \(A\) form an abelian
category, and finite presentations and their kernels remain coherent.
This follows by taking presentations over the coherent ring sheaf
\(\mathcal O_W\), using the coherent quotient
\(\mathcal O_W/\mathcal I_A\), and restricting to \(A\); the bridge
develops this argument in
\href{../complex-analytic-spaces-and-coherent-sheaves/coherent-sheaves-and-okas-theorem.html\#3-the-category-of-coherent-analytic-sheaves}{Coherent
sheaves and Oka's theorem, Section 3}.

\subsection{2. Local analytic algebra}\label{local-analytic-algebra}

\textbf{Preparation and division.} Write \(z=(z',w)\). Suppose the
holomorphic germ \(f\) satisfies \(f(0,w)=w^r v(w)\), with \(v(0)\ne0\).
Then, uniquely,

\[
f=uP,\qquad
P=w^r+a_1(z')w^{r-1}+\cdots+a_r(z'),\qquad a_j(0)=0,
\]

where \(u\) is an invertible holomorphic germ. Every holomorphic germ
\(g\) has a unique expression \(g=fq+R\), where \(R\) is a polynomial in
\(w\) of degree less than \(r\), with holomorphic coefficients in
\(z'\).

\textbf{Source and reuse.} This complete germ-level proof adapts
Demailly, II, Theorems 2.1 and 2.3, pp.~79--81, which credit C. L.
Siegel's contour argument. It is distributed under Demailly's
\href{https://www-fourier.univ-grenoble-alpes.fr/~demailly/documents.html}{explicit
OpenContent permission}, with that authorship retained. GPT-6.1 Sol at
Ultra effort reorganized the argument, made root multiplicities and germ
uniqueness explicit, and checked its use in Lemma 2.1. Demailly also
proves uniform bounded-function estimates; those stronger estimates are
not needed here. This component is not included in the CC0 dedication of
the independent lesson material.

\textbf{Proof.} The case \(r=0\) is division by a unit, with \(P=1\) and
\(R=0\). For \(r>0\), choose a small circle \(|\zeta|=\rho\) on which
\(f(0,\zeta)\ne0\), enclosing no zero other than zero. After shrinking
the \(z'\)-polydisc, \(f(z',\zeta)\) remains nonzero in an annulus
around that circle. Define

\[
s_k(z')=\frac{1}{2\pi i}\int_{|\zeta|=\rho}
\zeta^k\frac{\partial f/\partial w(z',\zeta)}{f(z',\zeta)}\,d\zeta.
\]

The scalar tools used here have complete earlier proofs in
\href{../foundations-of-von-neumann-algebras/cauchy-s-theorem-for-cycles-and-its-consequences.html}{Cauchy's
theorem for cycles and its consequences}: Lemma 0.1 proves interchange
of continuous integrals, Theorems 2.2--2.3 prove the disc integral
theorem and formula, and Theorems 3.2, 3.4 and 3.7 prove Taylor
expansion, Morera's theorem and isolatedness of zeros. These results
concern scalar holomorphic functions and do not require an
operator-algebra theorem.

The integrand is jointly continuous on the parameter polydisc times the
circle and holomorphic in each parameter coordinate. Its derivative
numerator is holomorphic by the polydisc Taylor theorem. Interchange a
triangle integral in one parameter coordinate with the circle integral,
using
\href{../foundations-of-von-neumann-algebras/cauchy-s-theorem-for-cycles-and-its-consequences.html\#OA-FND-CT-07}{Lemma
0.1(3)}; the triangle integral is zero by Cauchy's theorem.
\href{../foundations-of-von-neumann-algebras/cauchy-s-theorem-for-cycles-and-its-consequences.html\#OA-FND-CT-03}{Morera's
theorem, Theorem 3.4}, proves that each \(s_k\) is holomorphic in every
parameter coordinate. Joint continuity and the bridge's definition of
holomorphy make it holomorphic on the parameter polydisc.

For fixed \(z'\), the zeros in the closed disc are finite: they are
isolated by
\href{../foundations-of-von-neumann-algebras/cauchy-s-theorem-for-cycles-and-its-consequences.html\#OA-FND-CT-03}{Theorem
3.7}, none lies near the boundary, and an infinite subset of the
remaining compact disc would have an accumulation point. At a zero \(b\)
of multiplicity \(m_b\), Taylor factorization gives \(f=(w-b)^{m_b}h\),
with \(h(b)\ne0\), and therefore

\[
\frac{f_w}{f}=\frac{m_b}{w-b}+\frac{h'}h.
\]

Subtract \(\sum_b m_b/(w-b)\) from \(f_w/f\). The resulting function
extends holomorphically over every zero and throughout a slightly larger
disc. Its integral, also after multiplication by \(w^k\), is zero by
\href{../foundations-of-von-neumann-algebras/cauchy-s-theorem-for-cycles-and-its-consequences.html\#OA-FND-CT-02}{Theorem
2.2}. Cauchy's formula applied to the polynomial \(w^k\) gives

\[
s_k(z')=\sum_b m_b b^k\qquad(k\geq0).
\]

Here the term for \(k=0\) is \(m_b\), including when \(b=0\). Thus
\(s_0\) counts the zeros inside the circle, and \(s_k\) sums their
\(k\)-th powers with multiplicity. The integer \(s_0\) is continuous on
the connected parameter polydisc, hence constantly \(r\). Denote those
roots with repetition by \(b_1,\ldots,b_r\); individual roots need not
depend holomorphically on \(z'\).

Their elementary symmetric functions do: Newton's identities express
them recursively by \(c_0=1\) and

\[
k c_k=\sum_{j=1}^k(-1)^{j-1}c_{k-j}s_j.
\]

The polynomial \(P=\sum_{k=0}^r(-1)^k c_k w^{r-k}\) consequently has
holomorphic coefficients and exactly the same zeros, with
multiplicities, as \(f\) in the disc. At \(z'=0\), it is \(w^r\). For
each fixed \(z'\), both \(f/P\) and \(P/f\) have removable singularities
inside the circle. Their Cauchy integral representations using their
boundary values show that the extensions are jointly holomorphic in
\((z',w)\). They are inverse to one another. This proves preparation.
Any other monic polynomial of the stated form has the same small roots,
so equals \(P\); the unit is then determined.

For division by \(P\), choose a smaller circle inside a representative
of the germ \(g\), and shrink the \(z'\)-polydisc so all roots remain
inside it. Write its radius again as \(\rho\), and put

\[
Q(z',w)=\frac{1}{2\pi i}\int_{|\zeta|=\rho}
\frac{g(z',\zeta)}{P(z',\zeta)(\zeta-w)}\,d\zeta.
\]

Cauchy's formula gives

\[
g-PQ=\frac{1}{2\pi i}\int_{|\zeta|=\rho}
\frac{g(z',\zeta)}{P(z',\zeta)}
\frac{P(z',\zeta)-P(z',w)}{\zeta-w}\,d\zeta.
\]

The final divided difference is a polynomial in \(w\) of degree at most
\(r-1\), so this integral is the required \(R\). All coefficients and
\(Q\) are holomorphic. For \(f=uP\), take \(q=Q/u\). If two divisions
exist, their remainder difference is divisible by \(P\); it therefore
vanishes at its \(r\) roots with their multiplicities. A polynomial of
degree less than \(r\) with that property is zero. The quotient
difference is then zero. For germ uniqueness, shrink once more so all
roots lie in a neighbourhood on which both representations exist.
\(\square\)

\textbf{Lemma 2.1.} The ring \(H_n\) is Noetherian, and its maximal-adic
completion is \(\mathbf C[[z_1,\ldots,z_n]]\).

\textbf{Proof.} Induct on \(n\), beginning with \(H_0=\mathbf C\). For a
nonzero ideal \(J\), choose a nonzero \(f\in J\). A linear coordinate
change makes its first nonzero homogeneous part nonzero on the last
coordinate axis. Preparation then supplies a monic Weierstrass
polynomial \(P\in J\). Division identifies \(H_n/(P)\), as an
\(H_{n-1}\)-module, with the free module of polynomial remainders of
degrees \(0,\ldots,r-1\). By induction it is Noetherian. The image of
\(J\) is finitely generated in this quotient, so lifts of these
generators together with \(P\) generate \(J\). The zero ideal causes no
difficulty.

Modulo \((z)^m\), Taylor truncation identifies \(H_n/(z)^m\) with
polynomials of total degree less than \(m\). Taking the inverse limit
gives the full formal power-series ring. A convergent germ with every
Taylor coefficient zero is zero, so the Taylor map is injective.
\(\square\)

The same argument gives the completion of \(H_n/J\) as the formal
power-series ring modulo the extended ideal. Exactness of completion on
finite modules is proved using Artin--Rees in
\href{grothendiecks-existence-theorem.html\#2-exactness-and-the-actual-formal-kernel}{Grothendieck's
existence theorem, Section 2}, with the exact open completion results
cited there.

\textbf{Lemma 2.2 (holomorphic inverse and implicit functions).} A
holomorphic map between open subsets of \(\mathbf C^n\) whose derivative
is invertible at a point has a holomorphic inverse on smaller
neighbourhoods. If \(r\) equations have Jacobian rank \(r\), their zero
set is locally a complex manifold of dimension \(n-r\).

\textbf{Proof.} Translate the point and make a linear change so the map
is \(F(z)=z+g(z)\), where \(g(0)=0\) and \(dg(0)=0\). On a sufficiently
small closed ball, \(\|dg\|\leq c<1\). For \(w\) in a smaller ball, the
map \(z\mapsto w-g(z)\) maps the first ball into itself and is a
contraction. Its successive iterates converge uniformly, with a
geometric bound, to the unique solution \(z(w)\). The iterates are
holomorphic in \(w\), and local uniform convergence preserves
holomorphicity by Cauchy's formula. The estimate also gives injectivity
of \(F\) on the first ball. For the implicit statement, add
complementary coordinates to the given equations so that the combined
map has invertible derivative, and apply the inverse statement.
\(\square\)

Local analytic dimension is the Krull dimension of the local analytic
algebra.
\href{../complex-analytic-spaces-and-coherent-sheaves/analytic-germs-local-parametrization-and-the-nullstellensatz.html\#4-the-local-parametrization-theorem}{Local
finite parametrization, Theorem 4.1}, and
\href{../complex-analytic-spaces-and-coherent-sheaves/analytic-germs-local-parametrization-and-the-nullstellensatz.html\#5-the-nullstellensatz-and-dimension}{Dimension
of analytic germs, Theorem 5.4}, give the finite-projection proof: after
adapted linear coordinates, an irreducible germ is finite over a
polydisc of the dimension of its local ring and is a finite covering
away from a discriminant. This also treats dimension component by
component.

\subsection{3. Constructing the analytic space of a
variety}\label{constructing-the-analytic-space-of-a-variety}

For an affine variety \(X=\operatorname{Spec}(\mathbf C[z]/J)\), define
\(X^{\mathrm{an}}\) to have underlying set \(X(\mathbf C)\), its
ordinary topology as a closed subset of \(\mathbf C^n\), and the reduced
analytic structure defined above. A principal open \(D(h)\) is the
ordinary open subset where \(h\ne0\). Polynomial fractions with
denominator nonvanishing there are holomorphic.

We must verify that the analytic ideal is the extension of \(J\);
otherwise quotients could acquire an unnoticed reduction.

\textbf{Lemma 3.1.} For a radical algebraic ideal \(J\) and a complex
point \(x\), its extension to the ambient convergent local ring is
radical and equals the analytic vanishing ideal of its zero germ.

\textbf{Proof.} Translate \(x\) to zero and put

\[
R=\mathbf C[z]_{(z)},\qquad H=\mathbf C\{z\},\qquad S=\mathbf C[[z]].
\]

Both ambient completions are \(S\). The reduced local ring \(R/JR\) is
Nagata, and its completion \(S/JS\) is reduced. These are the precise
open algebraic proofs
\href{https://kokunoyumeto.github.io/stacks-zh-hans-cn/en/algebra.html\#algebra-proposition-nagata-universally-japanese}{Stacks,
Tags 0334 and 0335},
\href{https://kokunoyumeto.github.io/stacks-zh-hans-cn/en/algebra.html\#algebra-lemma-nagata-localize}{localization,
Tag 032U}, and
\href{https://kokunoyumeto.github.io/stacks-zh-hans-cn/en/more-algebra.html\#more-algebra-lemma-completion-reduced}{reduced
completion, Tag 07NZ}. The last proof reduces to the domain case
\href{https://kokunoyumeto.github.io/stacks-zh-hans-cn/en/algebra.html\#algebra-lemma-local-nagata-domain-analytically-unramified}{Tag
0331}.

The completion of \(H/JH\) is also \(S/JS\). Krull intersection embeds
the Noetherian local ring \(H/JH\) into this reduced completion, so it
is reduced. The analytic Nullstellensatz now identifies its vanishing
ideal with \(\sqrt{JH}=JH\). \(\square\)

\textbf{Theorem 3.2.} Every complex algebraic variety has an
analytification, uniquely characterized by the affine-chart
construction. Regular maps induce holomorphic maps, and analytification
respects composition, products, and reduced locally closed subvarieties.

\textbf{Proof.} On two affine charts, the algebraic transition maps are
regular maps and hence given locally by polynomial fractions. These are
holomorphic, as are their inverses. Lemma 3.1 identifies the
corresponding quotient structure sheaves. Thus the chart spaces glue
along precisely their ordinary open overlaps; the transition cocycle
gives a unique glued sheaf. Separateness makes the glued space
Hausdorff: near two points use affine charts, and the closed algebraic
diagonal becomes a closed analytic diagonal in their product. The same
chart argument glues a regular map and proves functoriality. Any other
space with these chart identifications glues by the same cocycle and
hence has a unique compatible isomorphism.

Product charts use the two disjoint sets of coordinate variables. Over
the perfect field \(\mathbf C\), a product of reduced varieties is
reduced, by the exact open proof
\href{https://kokunoyumeto.github.io/stacks-zh-hans-cn/en/varieties.html\#varieties-lemma-geometrically-reduced-any-base-change}{Stacks,
Tag 035Z}, so these are also the reduced analytic product charts. Closed
subvarieties follow from Lemma 3.1, and open subvarieties follow by
restriction. Combining the two proves the locally closed assertion.
\(\square\)

There is a morphism of locally ringed spaces \(a:X^{\mathrm{an}}\to X\),
taking a complex point to its closed scheme point and taking regular
functions to holomorphic functions. Zariski opens pull back to ordinary
opens, because their complements are algebraic zero sets. The ordinary
topology is finer: it contains small balls that need not be Zariski
open.

\subsection{4. The comparison at a
point}\label{the-comparison-at-a-point}

\textbf{Theorem 4.1.} For \(x\in X(\mathbf C)\), the local map

\[
A=\mathcal O_{X,x}\longrightarrow B=\mathcal O_{X^{\mathrm{an}},x}
\]

is faithfully flat, with residue field \(\mathbf C\), and induces an
isomorphism \(\widehat A\cong\widehat B\). Consequently the dimensions
and regularity of the two local rings agree.

\textbf{Proof.} On an ambient affine chart, Lemma 3.1 writes these rings
as the algebraic and convergent local quotients by the same polynomial
ideal. Taylor truncation and exact completion identify both completions
with \(S/JS\). Both rings are Noetherian. Their completion maps are
faithfully flat:
\href{https://kokunoyumeto.github.io/open-math-courses/courses/AG-CA/AG-CA-19.html\#3-noetherian-exactness-tensoring-and-flatness}{Completion,
Theorem 3.3} gives the proof using exact finite completion and the ideal
criterion for flatness. Thus \(\widehat B\) is flat over \(A\), and
faithfully flat over \(B\). For an injection of \(A\)-modules, its
tensor with \(B\) becomes injective after tensoring with \(\widehat B\);
faithful flatness over \(B\) detects that its kernel was zero. Therefore
\(B\) is flat over \(A\).

It is faithfully flat because it is a local flat map: a proper ideal of
\(A\) lies in its maximal ideal and cannot generate the unit ideal in
\(B\). Equivalently the usual ideal criterion for faithful flatness
applies. In particular \(IB\cap A=I\) for every ideal, and \(A\to B\) is
injective.

A Noetherian local ring and its completion have the same dimension and
the same cotangent space over the residue field; these are
\href{https://kokunoyumeto.github.io/open-math-courses/courses/AG-CA/AG-CA-19.html\#4-local-dimension-and-regularity}{Completion,
Theorem 4.1}, proved from the identical residue quotients and their
Hilbert--Samuel polynomial. The common completion therefore gives equal
dimensions and cotangent dimensions for \(A\) and \(B\). Regularity is
the equality of these two numbers, proving the assertion. \(\square\)

A Zariski dense open \(U\subset X\) is also ordinarily dense in
\(X^{\mathrm{an}}\). Indeed, if \(Z=X\setminus U\) contained an analytic
neighbourhood of \(x\), its analytic ideal at \(x\) would be zero.
Faithful flatness and Lemma 3.1 would give zero algebraic ideal there.
Coherence would then make \(Z\) contain a Zariski neighbourhood of
\(x\), contradicting density. A merely bijective continuous map,
however, need not preserve local rings; the cusp below shows why the
ring comparison matters.

\subsection{5. Coherent modules and
compactness}\label{coherent-modules-and-compactness}

Define

\[
\mathcal F^{\mathrm{an}}
=\mathcal O_{X^{\mathrm{an}}}\otimes_{a^{-1}\mathcal O_X}a^{-1}\mathcal F.
\]

\textbf{Proposition 5.1.} Analytification is exact and faithful on
coherent modules, preserves coherence, and commutes with tensor products
and internal Hom of coherent modules. For a closed immersion
\(i:Y\hookrightarrow X\), it commutes with \(i_*\).

\textbf{Proof.} The analytic stalk is \(B\otimes_A\mathcal F_x\), so
Theorem 4.1 proves exactness. A finite algebraic presentation becomes a
finite analytic presentation, whose cokernel is coherent by Section 1.
Tensor compatibility is associativity of tensor products.

For internal Hom, use a presentation \(A^r\to A^s\to M\to0\). Then
\(\operatorname{Hom}_A(M,N)\) is the kernel of \(N^s\to N^r\). Flatness
preserves that kernel, giving

\[
B\otimes_A\operatorname{Hom}_A(M,N)
\cong\operatorname{Hom}_B(B\otimes_A M,B\otimes_A N).
\]

The sheaf map is therefore an isomorphism on all analytic stalks. If the
analytification of a coherent module or a coherent-module morphism is
zero, faithful flatness makes its algebraic stalk or its image zero at
every closed point. A nonzero coherent module has nonempty closed
support, which on a finite-type complex scheme contains a closed point.
This proves faithfulness. Finally, closed-immersion pushforward regards
a module annihilated by the defining ideal as a module on the ambient
space. Tensoring gives exactly the same quotient-module description
analytically; stalks outside the closed set are zero. \(\square\)

\textbf{Proposition 5.2.} If \(X\) is projective over \(\mathbf C\),
then \(X^{\mathrm{an}}\) is compact.

\textbf{Proof.} The map from the unit sphere in \(\mathbf C^{n+1}\) to
complex projective space, sending a unit vector to its line, is
continuous and surjective. Thus \(\mathbf P^n(\mathbf C)\) is compact.
It is Hausdorff, for instance by identifying a line with its orthogonal
rank-one projection matrix. A closed algebraic subvariety is a closed
analytic subset by Theorem 3.2, hence compact. \(\square\)

The full properness criterion is \(X\) proper if and only if
\(X^{\mathrm{an}}\) is compact. Its
\href{https://kokunoyumeto.github.io/stacks-zh-hans-cn/en/gaga.html\#gaga-proposition-proper-if-and-only-if-compact-analytification}{open
proof in AI Integrated Stacks Project} uses Chow's lemma, assigned to
the planned \emph{Morphisms of schemes}, lesson \textbf{Projective
morphisms and Chow's lemma}. Concretely choose a proper surjection from
a quasi-projective variety \(V\) and close \(V\) in projective space. If
\(X\) is proper, \(V\) is already closed and its compact analytification
maps onto \(X^{\mathrm{an}}\). Conversely, compactness of
\(X^{\mathrm{an}}\) makes the closed graph's projection show that
\(V^{\mathrm{an}}\) is closed in its projective closure. It is dense by
Section 4, so \(V\) equals that closure; a proper surjective image of a
proper variety is proper.

\subsection{6. Three comparisons to keep
separate}\label{three-comparisons-to-keep-separate}

Affine space analytifies to \(\mathbf C^n\), with all holomorphic
functions on ordinary open subsets. The projective line analytifies to
the Riemann sphere, whose two standard coordinates are related by
\(w=1/z\).

For the cuspidal cubic \(y^2=x^3\), the parametrization
\(t\mapsto(t^2,t^3)\) is a continuous bijection
\(\mathbf C\to X^{\mathrm{an}}\). Away from the origin its inverse is
\(y/x\), and at the origin continuity follows from \(|t|^2=|x|\). It is
a homeomorphism. It is not an analytic isomorphism: at the cusp the
local ring is

\[
\mathbf C\{x,y\}/(y^2-x^3),
\]

of dimension one and cotangent dimension two. The relation has no linear
term. This agrees with the algebraic singularity by Theorem 4.1.
Moreover, substituting \((t^2,t^3)\) into a convergent germ never
produces a linear term in \(t\), so the inverse parameter is not
holomorphic at the cusp.

Finally, \(\exp:\mathbf C\to\mathbf C^*\) is holomorphic. An algebraic
morphism \(\mathbf A^1\to\mathbf G_m\) would send the invertible
coordinate of \(\mathbf G_m\) to a unit in \(\mathbf C[z]\). All such
units are nonzero constants. Thus the exponential is not algebraic.
Faithfulness of analytification does not imply fullness on nonproper
varieties.

\subsection{7. Exercises with solutions}\label{exercises-with-solutions}

\textbf{Exercise 7.1 (easy).} Identify the analytification of affine
space and explain why algebraic regular functions do not exhaust its
global analytic functions.

\textbf{Solution.} The ideal is zero in the chart construction, so the
underlying space is \(\mathbf C^n\) and its sheaf is
\(\mathcal O_{\mathbf C^n}\). For \(n\geq1\), the entire function
\(e^{z_1}\) is not polynomial: its derivatives of every order in \(z_1\)
are nonzero, whereas sufficiently high derivatives of a polynomial
vanish.

\textbf{Exercise 7.2 (easy).} Prove compactness of
\(\mathbf P^n(\mathbf C)\), including \(n=0\).

\textbf{Solution.} Every complex line meets the unit sphere. The
quotient map from that compact sphere is continuous and surjective, so
its image is compact. When \(n=0\), the quotient is one point. The
projection-matrix description in Proposition 5.2 also verifies
Hausdorffness.

\textbf{Exercise 7.3 (medium).} Compute the completions at zero of
\(\mathbf C[z]_{(z)}\) and \(\mathbf C\{z\}\), and the map between them.

\textbf{Solution.} Modulo \((z)^m\), both rings consist of the same
polynomials of total degree less than \(m\). In the algebraic
localization a denominator with nonzero constant term has a finite
geometric-series inverse modulo \((z)^m\). In the analytic ring Taylor
truncation gives the same inverse. Their compatible quotients therefore
have inverse limit \(\mathbf C[[z]]\), and the induced completion map is
the identity on every coefficient.

\textbf{Exercise 7.4 (medium).} Explain both exactness and detection of
zero for coherent modules under analytification.

\textbf{Solution.} At a complex point tensor with the faithfully flat
local algebra \(B\). Flatness preserves an exact sequence, so every
analytic stalk sequence is exact. Faithfulness implies that a finite
algebraic stalk with zero tensor is zero. If a coherent algebraic module
were nonzero, its closed support would contain a complex point,
contradicting this detection. Applying this to the coherent image of a
morphism proves detection of zero morphisms as well.

\textbf{Exercise 7.5 (medium).} Prove that no nonconstant algebraic
morphism \(\mathbf A^1\to\mathbf G_m\) exists. Why does this not
prohibit nonconstant holomorphic maps?

\textbf{Solution.} A ring map \(\mathbf C[u,u^{-1}]\to\mathbf C[z]\)
must send \(u\) to a polynomial unit, which is constant by degree
additivity in a product equal to one. Analytically, a nowhere-zero
entire function can be nonconstant; \(e^z\) and its reciprocal are both
entire. The analytic function ring has additional units.

\textbf{Exercise 7.6 (harder).} A homeomorphism need not remove a
singularity. Verify this on the cuspidal parametrization and locate the
failed inverse function.

\textbf{Solution.} The inverse is continuous off the cusp by the
quotient \(y/x\), and at the cusp by \(|t|=|x|^{1/2}\). The derivative
of the parametrization at zero is zero, so Lemma 2.2 does not apply.
More decisively, the analytic cusp ring embeds into \(\mathbf C\{t\}\)
by substitution, but its image has no term of degree one. Hence it
cannot contain the inverse coordinate \(t\). Its cotangent dimension is
two, whereas that of \(\mathbf C\{t\}\) is one; the homeomorphism is not
an isomorphism of locally ringed spaces.

\subsection{8. Proof providers and further
reading}\label{proof-providers-and-further-reading}

The definitions and comparison follow the mathematical framework of
J.-P. Serre, \emph{Géométrie algébrique et géométrie analytique},
Annales de l'Institut Fourier 6 (1956), 1--42. The article is credited
as the historical source; the proofs used here are the ones supplied in
this course and the specified openly reusable sources.

The
\href{https://kokunoyumeto.github.io/stacks-zh-hans-cn/en/gaga.html\#gaga-proposition-analytification}{AI
Integrated Stacks Project chapter on complex analytic spaces and GAGA}
gives arguments for the chart construction and local comparison,
including
\href{https://kokunoyumeto.github.io/stacks-zh-hans-cn/en/gaga.html\#gaga-theorem-analytification-completed-local-rings}{completed
local rings},
\href{https://kokunoyumeto.github.io/stacks-zh-hans-cn/en/gaga.html\#gaga-theorem-analytification-local-faithfully-flat}{faithful
flatness}, and
\href{https://kokunoyumeto.github.io/stacks-zh-hans-cn/en/gaga.html\#gaga-lemma-analytification-exact-faithful-coherent}{coherent
modules}. Its Oka--Cartan and Nullstellensatz inputs have programme
proof locations in Section 1. The incorporated preparation-and-division
proof retains Demailly's custom OpenContent permission and Siegel
attribution.

The statements proved in this lesson use the reduced-space convention
fixed at the start. Analytification of schemes with nilpotents requires
coherent quotient models by arbitrary analytic ideals, without replacing
them by radical ideals. That extension, including its chart gluing and
completed-local-ring comparison, remains a separate proof obligation;
coherent modules with nonreduced support do not establish it.

The algebraic completion prerequisite is supplied by \emph{Commutative
algebra for geometry}, lesson \textbf{Completion}, Theorems 3.1--3.3 and
4.1; Chow modification is planned in \emph{Morphisms of schemes}, lesson
\textbf{Projective morphisms and Chow's lemma}. Their exact open proof
providers are linked where used. The linked Stacks texts retain GFDL
1.2. Independently authored teaching material is CC0; the credited
preparation-and-division component retains Demailly's custom OpenContent
permission.
\end{document}
